\subsection{Лабораторная работа <<Расширенные сценарии работы с данными: отложенное получение данных, уведомления>> (1 неделя)}

\textbf{Цель работы:} Расширить приложение из предыдущей лабораторной работы механизмами обработки отсутствия сети и оповещения пользователя через уведомления.

\begin{itemize}
    \item При отсутствии сети пользовательский запрос сохраняется в историю со статусом «ожидает» и 
    ставится в очередь на фоновую загрузку через \textbf{WorkManager}.
    \item \texttt{Worker} запускается автоматически при появлении сети (через \texttt{Constraints}) 
    и использует существующий \texttt{Repository} из первой части работы.
    \item По завершении фоновой загрузки показывается \textbf{уведомление} с результатом (успех/ошибка) 
    и действием для перехода к данным.
    \item Реализуется экран \textbf{истории запросов} с отображением статуса каждого запроса 
    («в процессе», «успешно», «ошибка»).
\end{itemize}

\paragraph{Обязательные компоненты}

\begin{enumerate}
    \item \textbf{Единая таблица истории запросов} \\
    Расширьте существующую или создайте новую таблицу \texttt{search\_history} с полями:
    \begin{itemize}
        \item \texttt{id} — автоинкрементный идентификатор
        \item \texttt{query\_params} — параметры запроса (например, имя покемона, сериализовано в JSON)
        \item \texttt{timestamp} — время запроса
        \item \texttt{status} — статус: \texttt{PENDING} / \texttt{SUCCESS} / \texttt{FAILED}
        \item \texttt{result\_id} — идентификатор сохранённого результата в основной таблице кеша (для перехода)
    \end{itemize}
    
    \item \textbf{WorkManager для фоновой загрузки} \\
    Реализуйте \texttt{DataFetchWorker}, который:
    \begin{itemize}
        \item Выполняет запрос через существующий \texttt{Repository}
        \item При успехе: обновляет статус запроса на \texttt{SUCCESS}, сохраняет результат в кеш (как в первой части)
        \item При ошибке сети: завершается со статусом \texttt{Retry} (WorkManager применит бэк-офф автоматически)
        \item При 3 неудачах: обновляет статус на \texttt{FAILED}
        \item Показывает уведомление только при финальном результате (\texttt{SUCCESS} или \texttt{FAILED})
    \end{itemize}
    
    \item \textbf{Уведомления с действиями} \\
    Уведомления должны содержать:
    \begin{itemize}
        \item Заголовок: «Данные получены» / «Ошибка загрузки»
        \item Текст: краткое описание (например, «Покемон pikachu загружен»)
        \item Действие «Открыть» — \texttt{PendingIntent} для перехода на экран результата 
        (только при успехе)
        \item Действие «Повторить» — запуск нового \texttt{DataFetchWorker} с теми же параметрами 
        (только при ошибке)
    \end{itemize}
\end{enumerate}